\begin{abstract}
Kinase inhibitor selectivity is quantified using several competing definitions: 
S-score, selectivity entropy, Gini coefficient, and ratio-based measures,
which are applied interchangeably without justification.
Across four profiling datasets spanning three assay technologies (68-704
compounds each), we show these definitions cluster into two families measuring
different properties, and that rank instability is predictable from binding profile shape.
We trace instability to three sources: definitional noise below the detection
threshold, near-tied top targets, and distribution-ratio disagreement for
promiscuous compounds with a dominant target.
We derive four properties that any well-formed selectivity measure should
satisfy, then show that no existing definition satisfies all four. Finally, we propose a
measure that does, recovering its full-panel ranking from about a third of a panel.
\end{abstract}

\noindent\textbf{Keywords:} binding selectivity,
cancer drug discovery, cheminformatics, kinase inhibition, research methodology.
